% !TEX program = xelatex
% Bu dosya tex_uret.py tarafından sorular.py'den üretilmiştir; düzenleme için sorular.py'yi değiştirip yeniden üretin.
\documentclass[11pt,a4paper]{article}
\usepackage[margin=1.8cm,top=1.9cm,bottom=1.9cm]{geometry}
\usepackage{fontspec}
\usepackage{polyglossia}
\setdefaultlanguage{turkish}
\usepackage{amsmath,amssymb}
\usepackage{graphicx}
\usepackage{booktabs}
\usepackage{array}
\usepackage{enumitem}
\usepackage{multicol}
\usepackage[most]{tcolorbox}
\usepackage{fancyhdr}
\usepackage{titlesec}
\usepackage{xcolor}
\usepackage{needspace}
\usepackage[hidelinks]{hyperref}

\definecolor{anamavi}{RGB}{0,45,98}
\definecolor{cozumyesil}{RGB}{0,110,50}
\definecolor{uyarikirmizi}{RGB}{180,20,20}

\titleformat{\section}{\large\bfseries\color{anamavi}}{\thesection.}{0.5em}{}
\titleformat{\subsection}{\normalsize\bfseries\color{anamavi}}{\thesubsection}{0.5em}{}

\pagestyle{fancy}\fancyhf{}
\lhead{\small MMAT105 Genel Matematik}
\rhead{\small Hafta 15 -- Final Sınavı Deneme -- Çözümler}
\cfoot{\small\thepage}
\renewcommand{\headrulewidth}{0.4pt}

\newtcolorbox{kural}[1][Kural]{colback=anamavi!5,colframe=anamavi,title=#1,fonttitle=\bfseries,breakable,left=4pt,right=4pt,top=2pt,bottom=2pt}
\newtcolorbox{dikkat}[1][Dikkat]{colback=uyarikirmizi!5,colframe=uyarikirmizi,title=#1,fonttitle=\bfseries,breakable,left=4pt,right=4pt,top=2pt,bottom=2pt}
\newenvironment{cozum}{\begin{tcolorbox}[colback=cozumyesil!6,colframe=cozumyesil,boxrule=0.5pt,left=4pt,right=4pt,top=2pt,bottom=2pt]\textbf{\color{cozumyesil}Çözüm.}\ }{\end{tcolorbox}}
\newcommand{\sk}[1]{\textbf{\color{anamavi}#1)}~}
\setlength{\parindent}{0pt}
\setlength{\parskip}{3pt}
\setlength{\lineskiplimit}{2pt}\setlength{\lineskip}{4pt}
\setlength{\columnsep}{0.9cm}
\setlength{\columnseprule}{0.4pt}
\newlength{\soruara}\setlength{\soruara}{2.3cm}

\begin{document}

\begin{minipage}[c]{\textwidth}
{\LARGE\bfseries\color{anamavi} MMAT105 Genel Matematik}\\[4pt]
{\large\bfseries Final Sınavı Deneme -- Cevap Anahtarı ve Çözümler}\\[2pt]
{\normalsize Hafta 1--14 konuları}\\[6pt]
{\small 2026--2027 Güz}
\end{minipage}

\vspace{4pt}\noindent\textcolor{anamavi}{\rule{\textwidth}{1pt}}

\section*{Cevap Anahtarı}
\begin{center}
\renewcommand{\arraystretch}{1.5}
\begin{tabular}{|>{\centering\arraybackslash}p{1.1cm}|>{\centering\arraybackslash}p{1.1cm}|>{\centering\arraybackslash}p{1.1cm}|>{\centering\arraybackslash}p{1.1cm}|>{\centering\arraybackslash}p{1.1cm}|>{\centering\arraybackslash}p{1.1cm}|>{\centering\arraybackslash}p{1.1cm}|>{\centering\arraybackslash}p{1.1cm}|>{\centering\arraybackslash}p{1.1cm}|>{\centering\arraybackslash}p{1.1cm}|}\hline
\textbf{1} & \textbf{2} & \textbf{3} & \textbf{4} & \textbf{5} & \textbf{6} & \textbf{7} & \textbf{8} & \textbf{9} & \textbf{10}\\
\textbf{\color{cozumyesil}C} & \textbf{\color{cozumyesil}B} & \textbf{\color{cozumyesil}E} & \textbf{\color{cozumyesil}A} & \textbf{\color{cozumyesil}D} & \textbf{\color{cozumyesil}C} & \textbf{\color{cozumyesil}B} & \textbf{\color{cozumyesil}A} & \textbf{\color{cozumyesil}C} & \textbf{\color{cozumyesil}D}\\\hline
\textbf{11} & \textbf{12} & \textbf{13} & \textbf{14} & \textbf{15} & \textbf{16} & \textbf{17} & \textbf{18} & \textbf{19} & \textbf{20}\\
\textbf{\color{cozumyesil}B} & \textbf{\color{cozumyesil}A} & \textbf{\color{cozumyesil}E} & \textbf{\color{cozumyesil}E} & \textbf{\color{cozumyesil}B} & \textbf{\color{cozumyesil}D} & \textbf{\color{cozumyesil}D} & \textbf{\color{cozumyesil}A} & \textbf{\color{cozumyesil}E} & \textbf{\color{cozumyesil}C}\\\hline
\textbf{21} & \textbf{22} & \textbf{23} & \textbf{24} & \textbf{25} &  &  &  &  & \\
\textbf{\color{cozumyesil}B} & \textbf{\color{cozumyesil}E} & \textbf{\color{cozumyesil}C} & \textbf{\color{cozumyesil}A} & \textbf{\color{cozumyesil}D} &  &  &  &  & \\\hline
\end{tabular}
\end{center}
{\small\textbf{Satır hâlinde:} 1-C, \ 2-B, \ 3-E, \ 4-A, \ 5-D, \ 6-C, \ 7-B, \ 8-A, \ 9-C, \ 10-D, \ 11-B, \ 12-A, \ 13-E, \ 14-E, \ 15-B, \ 16-D, \ 17-D, \ 18-A, \ 19-E, \ 20-C, \ 21-B, \ 22-E, \ 23-C, \ 24-A, \ 25-D}


\section*{Soruların Konu ve Zorluk Dağılımı}
\begin{center}\small\setlength{\tabcolsep}{4pt}
\begin{tabular}{@{}r l l c@{\hspace{1.2em}}|@{\hspace{1.2em}}r l l c@{}}
\toprule
No & Hafta -- Konu & Düzey & Cevap & No & Hafta -- Konu & Düzey & Cevap\\
\midrule
1 & H1 -- İşlem önceliği, işaret & Kolay & C & 14 & H11 -- İki kare farkı & Kolay & E\\
2 & H2 -- Asal çarpanlar & Kolay & B & 15 & H11 -- Tam kare özdeşliği & Orta & B\\
3 & H3 -- EKOK problemi & Orta & E & 16 & H11 -- Rasyonel ifade sadeleştirme & Orta & D\\
4 & H4 -- Merdiven kesir & Orta & A & 17 & H12 -- Oran & Kolay & D\\
5 & H5 -- Denklem kurma & Orta & D & 18 & H12 -- Bileşik orantı & Orta & A\\
6 & H6 -- İki bilinmeyenli problem & Orta & C & 19 & H12 -- Aritmetik ortalama & Zor & E\\
7 & H7 -- Mutlak değerli eşitsizlik & Zor & B & 20 & H13 -- Sayı problemi & Kolay & C\\
8 & H9 -- Üslü sayılarda işlem & Kolay & A & 21 & H13 -- Yaş problemi & Orta & B\\
9 & H9 -- Ortak çarpan (üslü) & Orta & C & 22 & H13 -- Kesir problemi & Orta & E\\
10 & H9 -- Üslü sayılarda sıralama & Zor & D & 23 & H14 -- Basit faiz & Kolay & C\\
11 & H10 -- Kök dışına çıkarma & Kolay & B & 24 & H14 -- İndirim ve kâr & Orta & A\\
12 & H10 -- Paydayı rasyonel yapma & Orta & A & 25 & H14 -- Karışım & Zor & D\\
13 & H10 -- İç içe kök & Zor & E &  & & & \\
\bottomrule\end{tabular}
\end{center}
{\small Düzey dağılımı: Kolay 8, Orta 12, Zor (TYT orta üstü) 5. \quad Doğru şık dağılımı: A: 5, B: 5, C: 5, D: 5, E: 5.}

\needspace{8\baselineskip}
\section*{Çözümler}

\par\medskip\noindent\begin{minipage}{\linewidth}
\noindent\textbf{\color{anamavi}Soru 1.} {\small\color{gray}(Hafta 1 -- İşlem önceliği, işaret; Kolay)}\par
\((-2)^3-(-3)^2+4\cdot 5-12\div 3\) işleminin sonucu kaçtır?
\par\smallskip\noindent{\setlength{\tabcolsep}{1pt}\begin{tabular*}{\linewidth}{@{}l@{\extracolsep{\fill}}llll@{}}\sk{A}\(-17\) & \sk{B}\(-9\) & \sk{C}\(-1\) & \sk{D}\(1\) & \sk{E}\(17\)\end{tabular*}}
\begin{cozum}
Üsler: \((-2)^3=-8\), \((-3)^2=9\). Çarpma ve bölme: \(4\cdot 5=20\), \(12\div 3=4\).
\[-8-9+20-4=-1.\]
\textit{Tipik hata:} \(-(-3)^2\) ifadesini \(+9\) almak \(\Rightarrow 17\).
\par\smallskip\noindent\textbf{\color{cozumyesil}Cevap: C}
\end{cozum}
\end{minipage}

\par\medskip\noindent\begin{minipage}{\linewidth}
\noindent\textbf{\color{anamavi}Soru 2.} {\small\color{gray}(Hafta 2 -- Asal çarpanlar; Kolay)}\par
\(84\) sayısının farklı asal bölenlerinin toplamı kaçtır?
\par\smallskip\noindent{\setlength{\tabcolsep}{1pt}\begin{tabular*}{\linewidth}{@{}l@{\extracolsep{\fill}}llll@{}}\sk{A}\(7\) & \sk{B}\(12\) & \sk{C}\(14\) & \sk{D}\(16\) & \sk{E}\(21\)\end{tabular*}}
\begin{cozum}
\(84=2^2\cdot 3\cdot 7\). Farklı asal bölenler \(2,\ 3,\ 7\): \ \(2+3+7=12\).
\textit{Tipik hata:} \(2\)'yi iki kez saymak (\(2+2+3+7=14\)).
\par\smallskip\noindent\textbf{\color{cozumyesil}Cevap: B}
\end{cozum}
\end{minipage}

\par\medskip\noindent\begin{minipage}{\linewidth}
\noindent\textbf{\color{anamavi}Soru 3.} {\small\color{gray}(Hafta 3 -- EKOK problemi; Orta)}\par
Bir duraktan üç otobüs hattının otobüsleri sırasıyla \(15\), \(20\) ve \(25\) dakikada bir kalkmaktadır. Üç hattın otobüsleri saat 08.00'de birlikte kalktığına göre, üçü yeniden ilk kez birlikte saat kaçta kalkar?
\par\smallskip\noindent{\setlength{\tabcolsep}{1pt}\begin{tabular*}{\linewidth}{@{}l@{\extracolsep{\fill}}llll@{}}\sk{A}09.00 & \sk{B}09.15 & \sk{C}09.40 & \sk{D}12.00 & \sk{E}13.00\end{tabular*}}
\begin{cozum}
Birlikte kalkış aralığı EKOK\((15,20,25)\)'tir:
\(15=3\cdot 5,\ 20=2^2\cdot 5,\ 25=5^2\Rightarrow\) EKOK \(=2^2\cdot 3\cdot 5^2=300\) dakika \(=5\) saat.
\(08.00+5\) saat \(=\) \textbf{13.00}.
\textit{Tipik hatalar:} yalnızca iki hattın EKOK'unu almak: EKOK\((15,20)=60\) dk \(\Rightarrow\) 09.00; EKOK\((20,25)=100\) dk \(\Rightarrow\) 09.40; EKOK\((15,25)=75\) dk \(\Rightarrow\) 09.15.
\par\smallskip\noindent\textbf{\color{cozumyesil}Cevap: E}
\end{cozum}
\end{minipage}

\par\medskip\noindent\begin{minipage}{\linewidth}
\noindent\textbf{\color{anamavi}Soru 4.} {\small\color{gray}(Hafta 4 -- Merdiven kesir; Orta)}\par
\(2-\cfrac{1}{1-\cfrac{1}{3}}\) işleminin sonucu kaçtır?
\par\smallskip\noindent{\setlength{\tabcolsep}{1pt}\begin{tabular*}{\linewidth}{@{}l@{\extracolsep{\fill}}llll@{}}\sk{A}\(\dfrac{1}{2}\) & \sk{B}\(\dfrac{2}{3}\) & \sk{C}\(\dfrac{4}{3}\) & \sk{D}\(\dfrac{3}{2}\) & \sk{E}\(\dfrac{5}{2}\)\end{tabular*}}
\begin{cozum}
En alttan başlayalım: \(1-\dfrac{1}{3}=\dfrac{2}{3}\). \quad \(\dfrac{1}{2/3}=\dfrac{3}{2}\).
\[2-\frac{3}{2}=\frac{1}{2}.\]
\textit{Tipik hata:} \(\dfrac{1}{2/3}\)'ü \(\dfrac{2}{3}\) almak \(\Rightarrow 2-\dfrac{2}{3}=\dfrac{4}{3}\).
\par\smallskip\noindent\textbf{\color{cozumyesil}Cevap: A}
\end{cozum}
\end{minipage}

\par\medskip\noindent\begin{minipage}{\linewidth}
\noindent\textbf{\color{anamavi}Soru 5.} {\small\color{gray}(Hafta 5 -- Denklem kurma; Orta)}\par
Ali'nin parası Ayşe'nin parasının \(3\) katıdır. Ali, Ayşe'ye \(40\) TL verirse paraları eşit oluyor. Ali'nin başlangıçtaki parası kaç TL'dir?
\par\smallskip\noindent{\setlength{\tabcolsep}{1pt}\begin{tabular*}{\linewidth}{@{}l@{\extracolsep{\fill}}llll@{}}\sk{A}\(40\) & \sk{B}\(80\) & \sk{C}\(100\) & \sk{D}\(120\) & \sk{E}\(160\)\end{tabular*}}
\begin{cozum}
Ayşe'nin parası \(x\) ise Ali'ninki \(3x\). Ali \(40\) TL verince Ali'de \(3x-40\), Ayşe'de \(x+40\) olur:
\[3x-40=x+40\Rightarrow 2x=80\Rightarrow x=40.\]
Ali'nin başlangıçtaki parası \(3\cdot 40=120\) TL. (Kontrol: \(80=80\) \checkmark)
\textit{Tipik hatalar:} \(x=40\)'ı (Ayşe'nin parası) cevap sanmak; toplamı (\(160\)) vermek.
\par\smallskip\noindent\textbf{\color{cozumyesil}Cevap: D}
\end{cozum}
\end{minipage}

\par\medskip\noindent\begin{minipage}{\linewidth}
\noindent\textbf{\color{anamavi}Soru 6.} {\small\color{gray}(Hafta 6 -- İki bilinmeyenli problem; Orta)}\par
\(30\) soruluk bir sınavda her doğru cevap \(4\) puan kazandırmakta, her yanlış cevap \(1\) puan kaybettirmektedir. Tüm soruları cevaplayan bir öğrenci \(70\) puan aldığına göre kaç soruya doğru cevap vermiştir?
\par\smallskip\noindent{\setlength{\tabcolsep}{1pt}\begin{tabular*}{\linewidth}{@{}l@{\extracolsep{\fill}}llll@{}}\sk{A}\(14\) & \sk{B}\(18\) & \sk{C}\(20\) & \sk{D}\(22\) & \sk{E}\(25\)\end{tabular*}}
\begin{cozum}
Doğru sayısı \(d\), yanlış sayısı \(y\): \ \(d+y=30\) ve \(4d-y=70\).
Taraf tarafa toplayalım: \(5d=100\Rightarrow d=20\), \(y=10\). (Kontrol: \(80-10=70\) \checkmark)
\par\smallskip\noindent\textbf{\color{cozumyesil}Cevap: C}
\end{cozum}
\end{minipage}

\par\medskip\noindent\begin{minipage}{\linewidth}
\noindent\textbf{\color{anamavi}Soru 7.} {\small\color{gray}(Hafta 7 -- Mutlak değerli eşitsizlik; Zor)}\par
\(2<|x-1|\le 5\) eşitsizliğini sağlayan kaç tane \(x\) tam sayısı vardır?
\par\smallskip\noindent{\setlength{\tabcolsep}{1pt}\begin{tabular*}{\linewidth}{@{}l@{\extracolsep{\fill}}llll@{}}\sk{A}\(4\) & \sk{B}\(6\) & \sk{C}\(7\) & \sk{D}\(8\) & \sk{E}\(10\)\end{tabular*}}
\begin{cozum}
\(|x-1|\le 5\iff -4\le x\le 6\). \quad \(|x-1|>2\iff x-1>2\) veya \(x-1<-2\iff x>3\) veya \(x<-1\).
Ortak çözüm: \(-4\le x<-1\) veya \(3<x\le 6\).
Tam sayılar: \(-4,-3,-2\) ve \(4,5,6\) \(\Rightarrow\) \(6\) tane.
\textit{Tipik hata:} \(2\)'yi dahil etmek (\(|x-1|\ge 2\)) \(\Rightarrow\) \(-1\) ve \(3\) de eklenir: \(8\).
\par\smallskip\noindent\textbf{\color{cozumyesil}Cevap: B}
\end{cozum}
\end{minipage}

\par\medskip\noindent\begin{minipage}{\linewidth}
\noindent\textbf{\color{anamavi}Soru 8.} {\small\color{gray}(Hafta 9 -- Üslü sayılarda işlem; Kolay)}\par
\(\dfrac{2^5\cdot 4^3}{8^4}\) işleminin sonucu kaçtır?
\par\smallskip\noindent{\setlength{\tabcolsep}{1pt}\begin{tabular*}{\linewidth}{@{}l@{\extracolsep{\fill}}llll@{}}\sk{A}\(\dfrac{1}{2}\) & \sk{B}\(2\) & \sk{C}\(4\) & \sk{D}\(8\) & \sk{E}\(16\)\end{tabular*}}
\begin{cozum}
Hepsini \(2\) tabanına çevirelim: \(4^3=(2^2)^3=2^6\), \(8^4=(2^3)^4=2^{12}\).
\[\frac{2^5\cdot 2^6}{2^{12}}=2^{5+6-12}=2^{-1}=\frac{1}{2}.\]
\textit{Tipik hata:} \((2^3)^4\) için üsleri toplamak (\(2^7\)) \(\Rightarrow 2^4=16\).
\par\smallskip\noindent\textbf{\color{cozumyesil}Cevap: A}
\end{cozum}
\end{minipage}

\par\medskip\noindent\begin{minipage}{\linewidth}
\noindent\textbf{\color{anamavi}Soru 9.} {\small\color{gray}(Hafta 9 -- Ortak çarpan (üslü); Orta)}\par
\(x\) bir gerçek sayı olmak üzere \(\dfrac{2^{x+4}-2^{x}}{2^{x+1}+2^{x}}\) ifadesinin değeri kaçtır?
\par\smallskip\noindent{\setlength{\tabcolsep}{1pt}\begin{tabular*}{\linewidth}{@{}l@{\extracolsep{\fill}}llll@{}}\sk{A}\(3\) & \sk{B}\(4\) & \sk{C}\(5\) & \sk{D}\(\dfrac{16}{3}\) & \sk{E}\(8\)\end{tabular*}}
\begin{cozum}
Pay ve paydada \(2^x\) ortak çarpan parantezine alınır:
\[\frac{2^x(2^4-1)}{2^x(2+1)}=\frac{16-1}{2+1}=\frac{15}{3}=5.\]
\textit{Tipik hatalar:} \(2^{x+4}-2^x=2^4\) gibi üsleri çıkarmak \(\Rightarrow\) \(\dfrac{16}{3}\); paydada \(2^{x+1}+2^x=2^{2x+1}\) sanıp \(2^x\)'leri yanlış sadeleştirmek.
\par\smallskip\noindent\textbf{\color{cozumyesil}Cevap: C}
\end{cozum}
\end{minipage}

\par\medskip\noindent\begin{minipage}{\linewidth}
\noindent\textbf{\color{anamavi}Soru 10.} {\small\color{gray}(Hafta 9 -- Üslü sayılarda sıralama; Zor)}\par
\(a=2^{48},\quad b=3^{36},\quad c=5^{24}\) sayılarının doğru sıralanışı aşağıdakilerden hangisidir?
\par\smallskip\noindent\begin{tabular}{@{}p{0.47\linewidth}p{0.47\linewidth}@{}}\sk{A}\(a<b<c\) & \sk{B}\(b<c<a\)\\[7pt] \sk{C}\(c<a<b\) & \sk{D}\(a<c<b\)\\[7pt] \sk{E}\(c<b<a\) & \end{tabular}
\begin{cozum}
Üslerin EBOB'u \(12\): hepsini \(12\). kuvvet olarak yazalım.
\[a=(2^4)^{12}=16^{12},\qquad b=(3^3)^{12}=27^{12},\qquad c=(5^2)^{12}=25^{12}.\]
Üsler eşit ve tabanlar pozitif: \(16<25<27\Rightarrow a<c<b\).
\textit{Tipik hata:} yalnızca tabanlara (\(2<3<5\)) ya da yalnızca üslere bakmak.
\par\smallskip\noindent\textbf{\color{cozumyesil}Cevap: D}
\end{cozum}
\end{minipage}

\par\medskip\noindent\begin{minipage}{\linewidth}
\noindent\textbf{\color{anamavi}Soru 11.} {\small\color{gray}(Hafta 10 -- Kök dışına çıkarma; Kolay)}\par
\(\sqrt{98}+\sqrt{8}-\sqrt{32}\) işleminin sonucu kaçtır?
\par\smallskip\noindent\begin{tabular}{@{}p{0.47\linewidth}p{0.47\linewidth}@{}}\sk{A}\(\sqrt{74}\) & \sk{B}\(5\sqrt{2}\)\\[7pt] \sk{C}\(9\sqrt{2}\) & \sk{D}\(13\sqrt{2}\)\\[7pt] \sk{E}\(3\sqrt{2}\) & \end{tabular}
\begin{cozum}
\(\sqrt{98}=\sqrt{49\cdot 2}=7\sqrt2\), \ \(\sqrt{8}=2\sqrt2\), \ \(\sqrt{32}=\sqrt{16\cdot 2}=4\sqrt2\).
\[7\sqrt2+2\sqrt2-4\sqrt2=5\sqrt2.\]
\textit{Tipik hatalar:} kök içlerini toplamak: \(\sqrt{98+8-32}=\sqrt{74}\) (yanlış!); \(-\sqrt{32}\)'nin işaretini kaçırmak \(\Rightarrow 13\sqrt2\).
\par\smallskip\noindent\textbf{\color{cozumyesil}Cevap: B}
\end{cozum}
\end{minipage}

\par\medskip\noindent\begin{minipage}{\linewidth}
\noindent\textbf{\color{anamavi}Soru 12.} {\small\color{gray}(Hafta 10 -- Paydayı rasyonel yapma; Orta)}\par
\(\dfrac{6}{\sqrt{7}-1}-\sqrt{7}\) işleminin sonucu kaçtır?
\par\smallskip\noindent{\setlength{\tabcolsep}{1pt}\begin{tabular*}{\linewidth}{@{}l@{\extracolsep{\fill}}llll@{}}\sk{A}\(1\) & \sk{B}\(\sqrt{7}-1\) & \sk{C}\(2\) & \sk{D}\(\sqrt{7}+1\) & \sk{E}\(6\)\end{tabular*}}
\begin{cozum}
Eşlenikle genişletelim:
\[\frac{6}{\sqrt7-1}\cdot\frac{\sqrt7+1}{\sqrt7+1}=\frac{6(\sqrt7+1)}{7-1}=\sqrt7+1.\]
\((\sqrt7+1)-\sqrt7=1\).
\textit{Tipik hata:} \(\sqrt7\)'yi çıkarmayı unutmak \(\Rightarrow \sqrt7+1\).
\par\smallskip\noindent\textbf{\color{cozumyesil}Cevap: A}
\end{cozum}
\end{minipage}

\par\medskip\noindent\begin{minipage}{\linewidth}
\noindent\textbf{\color{anamavi}Soru 13.} {\small\color{gray}(Hafta 10 -- İç içe kök; Zor)}\par
\(\sqrt{11+\sqrt{72}}\) ifadesinin eşiti aşağıdakilerden hangisidir?
\par\smallskip\noindent\begin{tabular}{@{}p{0.47\linewidth}p{0.47\linewidth}@{}}\sk{A}\(\sqrt{6}+\sqrt{5}\) & \sk{B}\(2\sqrt{2}+\sqrt{3}\)\\[7pt] \sk{C}\(2+\sqrt{7}\) & \sk{D}\(3+2\sqrt{2}\)\\[7pt] \sk{E}\(3+\sqrt{2}\) & \end{tabular}
\begin{cozum}
Kural: \(\sqrt{a+b+2\sqrt{ab}}=\sqrt a+\sqrt b\) \ (\(a,b>0\)). Önce iç kökü \(2\sqrt{\ \ }\) biçimine getirelim:
\(\sqrt{72}=\sqrt{4\cdot 18}=2\sqrt{18}\). Toplamı \(11\), çarpımı \(18\) olan sayılar: \(9\) ve \(2\).
\[\sqrt{11+2\sqrt{18}}=\sqrt9+\sqrt2=3+\sqrt2.\]
Kontrol: \((3+\sqrt2)^2=9+6\sqrt2+2=11+\sqrt{72}\) \checkmark \
\textit{Tipik hata:} \(2\) çarpanını ayırmadan toplamı \(11\) olan sayılar aramak (\(6\) ve \(5\)).
\par\smallskip\noindent\textbf{\color{cozumyesil}Cevap: E}
\end{cozum}
\end{minipage}

\par\medskip\noindent\begin{minipage}{\linewidth}
\noindent\textbf{\color{anamavi}Soru 14.} {\small\color{gray}(Hafta 11 -- İki kare farkı; Kolay)}\par
\(x^2-y^2=48\) ve \(x+y=12\) olduğuna göre \(x\cdot y\) kaçtır?
\par\smallskip\noindent{\setlength{\tabcolsep}{1pt}\begin{tabular*}{\linewidth}{@{}l@{\extracolsep{\fill}}llll@{}}\sk{A}\(4\) & \sk{B}\(16\) & \sk{C}\(20\) & \sk{D}\(24\) & \sk{E}\(32\)\end{tabular*}}
\begin{cozum}
\(x^2-y^2=(x-y)(x+y)\Rightarrow 48=(x-y)\cdot 12\Rightarrow x-y=4\).
\(x+y=12\) ile toplarsak \(2x=16\Rightarrow x=8\), \(y=4\). \quad \(x\cdot y=32\).
\textit{Tipik hata:} \(x-y=4\)'ü cevap sanmak.
\par\smallskip\noindent\textbf{\color{cozumyesil}Cevap: E}
\end{cozum}
\end{minipage}

\par\medskip\noindent\begin{minipage}{\linewidth}
\noindent\textbf{\color{anamavi}Soru 15.} {\small\color{gray}(Hafta 11 -- Tam kare özdeşliği; Orta)}\par
\(a+b=10\) ve \(a\cdot b=21\) olduğuna göre \(a^2+b^2\) kaçtır?
\par\smallskip\noindent{\setlength{\tabcolsep}{1pt}\begin{tabular*}{\linewidth}{@{}l@{\extracolsep{\fill}}llll@{}}\sk{A}\(16\) & \sk{B}\(58\) & \sk{C}\(79\) & \sk{D}\(100\) & \sk{E}\(142\)\end{tabular*}}
\begin{cozum}
\((a+b)^2=a^2+2ab+b^2\Rightarrow a^2+b^2=(a+b)^2-2ab=100-42=58\).
\textit{Tipik hatalar:} \(2ab\) yerine \(ab\) çıkarmak \(\Rightarrow 79\); \((a+b)^2=a^2+b^2\) sanmak \(\Rightarrow 100\); \(2ab\)'yi eklemek \(\Rightarrow 142\).
\par\smallskip\noindent\textbf{\color{cozumyesil}Cevap: B}
\end{cozum}
\end{minipage}

\par\medskip\noindent\begin{minipage}{\linewidth}
\noindent\textbf{\color{anamavi}Soru 16.} {\small\color{gray}(Hafta 11 -- Rasyonel ifade sadeleştirme; Orta)}\par
\(x\notin\{-3,\,-2,\,0,\,3\}\) olmak üzere \(\dfrac{x^2-x-6}{x^2-9}\div\dfrac{x+2}{x}\) ifadesinin en sade biçimi aşağıdakilerden hangisidir?
\par\smallskip\noindent\begin{tabular}{@{}p{0.47\linewidth}p{0.47\linewidth}@{}}\sk{A}\(\dfrac{x+2}{x+3}\) & \sk{B}\(\dfrac{x}{x-3}\)\\[7pt] \sk{C}\(\dfrac{x-3}{x}\) & \sk{D}\(\dfrac{x}{x+3}\)\\[7pt] \sk{E}\(\dfrac{(x+2)^2}{x(x+3)}\) & \end{tabular}
\begin{cozum}
\(x^2-x-6=(x-3)(x+2)\) (çarpımı \(-6\), toplamı \(-1\) olan sayılar: \(-3\) ve \(2\)); \ \(x^2-9=(x-3)(x+3)\).
Bölme, ikinci kesrin ters çevrilip çarpılmasıdır:
\[\frac{(x-3)(x+2)}{(x-3)(x+3)}\cdot\frac{x}{x+2}=\frac{x}{x+3}.\]
\textit{Tipik hata:} ters çevirmeden çarpmak \(\Rightarrow \dfrac{(x+2)^2}{x(x+3)}\).
\par\smallskip\noindent\textbf{\color{cozumyesil}Cevap: D}
\end{cozum}
\end{minipage}

\par\medskip\noindent\begin{minipage}{\linewidth}
\noindent\textbf{\color{anamavi}Soru 17.} {\small\color{gray}(Hafta 12 -- Oran; Kolay)}\par
\(\dfrac{a}{b}=\dfrac{5}{8}\) ve \(a+b=91\) olduğuna göre \(b\) kaçtır?
\par\smallskip\noindent{\setlength{\tabcolsep}{1pt}\begin{tabular*}{\linewidth}{@{}l@{\extracolsep{\fill}}llll@{}}\sk{A}\(21\) & \sk{B}\(35\) & \sk{C}\(49\) & \sk{D}\(56\) & \sk{E}\(63\)\end{tabular*}}
\begin{cozum}
\(a=5k,\ b=8k\Rightarrow 13k=91\Rightarrow k=7\). \quad \(b=8\cdot 7=56\) \ (\(a=35\)).
\textit{Tipik hatalar:} \(a\)'yı (\(35\)) ya da \(b-a\)'yı (\(21\)) vermek.
\par\smallskip\noindent\textbf{\color{cozumyesil}Cevap: D}
\end{cozum}
\end{minipage}

\par\medskip\noindent\begin{minipage}{\linewidth}
\noindent\textbf{\color{anamavi}Soru 18.} {\small\color{gray}(Hafta 12 -- Bileşik orantı; Orta)}\par
\(8\) bahçıvan günde \(6\) saat çalışarak bir parkın bakımını \(5\) günde bitirmektedir. Aynı hızda çalışan \(6\) bahçıvan günde \(4\) saat çalışarak aynı işi kaç günde bitirir?
\par\smallskip\noindent{\setlength{\tabcolsep}{1pt}\begin{tabular*}{\linewidth}{@{}l@{\extracolsep{\fill}}llll@{}}\sk{A}\(10\) & \sk{B}\(12\) & \sk{C}\(15\) & \sk{D}\(16\) & \sk{E}\(20\)\end{tabular*}}
\begin{cozum}
İş miktarı \(=\) kişi \(\times\) günlük saat \(\times\) gün sabittir:
\[8\cdot 6\cdot 5=6\cdot 4\cdot g\Rightarrow 240=24g\Rightarrow g=10.\]
\textit{Not:} Bahçıvan sayısı ve günlük çalışma süresi azaldığından gün sayısı artar (ters orantı); oranları ters kurmak en sık hatadır.
\par\smallskip\noindent\textbf{\color{cozumyesil}Cevap: A}
\end{cozum}
\end{minipage}

\par\medskip\noindent\begin{minipage}{\linewidth}
\noindent\textbf{\color{anamavi}Soru 19.} {\small\color{gray}(Hafta 12 -- Aritmetik ortalama; Zor)}\par
Bir sınıftaki \(24\) öğrencinin not ortalaması \(65\)'tir. Sınıfa, not ortalaması \(80\) olan \(k\) öğrenci katıldığında sınıfın not ortalaması \(70\) oluyor. Buna göre \(k\) kaçtır?
\par\smallskip\noindent{\setlength{\tabcolsep}{1pt}\begin{tabular*}{\linewidth}{@{}l@{\extracolsep{\fill}}llll@{}}\sk{A}\(4\) & \sk{B}\(6\) & \sk{C}\(8\) & \sk{D}\(10\) & \sk{E}\(12\)\end{tabular*}}
\begin{cozum}
Toplamlar üzerinden: \(24\cdot 65+80k=70(24+k)\).
\[1560+80k=1680+70k\Rightarrow 10k=120\Rightarrow k=12.\]
(Kısa yol: ilk grup ortalamanın \(5\) altında, yeni grup \(10\) üstünde: \(24\cdot 5=k\cdot 10\Rightarrow k=12\).)
\par\smallskip\noindent\textbf{\color{cozumyesil}Cevap: E}
\end{cozum}
\end{minipage}

\par\medskip\noindent\begin{minipage}{\linewidth}
\noindent\textbf{\color{anamavi}Soru 20.} {\small\color{gray}(Hafta 13 -- Sayı problemi; Kolay)}\par
Hangi sayının \(\dfrac{2}{5}\)'inin \(3\) fazlası \(15\)'tir?
\par\smallskip\noindent{\setlength{\tabcolsep}{1pt}\begin{tabular*}{\linewidth}{@{}l@{\extracolsep{\fill}}llll@{}}\sk{A}\(18\) & \sk{B}\(24\) & \sk{C}\(30\) & \sk{D}\(36\) & \sk{E}\(45\)\end{tabular*}}
\begin{cozum}
\(\dfrac{2}{5}x+3=15\Rightarrow\dfrac{2}{5}x=12\Rightarrow x=12\cdot\dfrac{5}{2}=30\).
\textit{Tipik hata:} \(3\)'ü çıkarmak yerine eklemek: \(\dfrac25x=18\Rightarrow x=45\).
\par\smallskip\noindent\textbf{\color{cozumyesil}Cevap: C}
\end{cozum}
\end{minipage}

\par\medskip\noindent\begin{minipage}{\linewidth}
\noindent\textbf{\color{anamavi}Soru 21.} {\small\color{gray}(Hafta 13 -- Yaş problemi; Orta)}\par
Bir annenin yaşı ile kızının yaşının toplamı bugün \(50\)'dir. \(4\) yıl önce annenin yaşı, kızının yaşının \(5\) katıydı. Annenin bugünkü yaşı kaçtır?
\par\smallskip\noindent{\setlength{\tabcolsep}{1pt}\begin{tabular*}{\linewidth}{@{}l@{\extracolsep{\fill}}llll@{}}\sk{A}\(35\) & \sk{B}\(39\) & \sk{C}\(41\) & \sk{D}\(45\) & \sk{E}\(46\)\end{tabular*}}
\begin{cozum}
Annenin bugünkü yaşı \(a\) ise kızınki \(50-a\). \(4\) yıl önce: \(a-4=5\,(50-a-4)\).
\[a-4=230-5a\Rightarrow 6a=234\Rightarrow a=39.\]
Kontrol: kız \(11\) yaşında; \(4\) yıl önce \(35=5\cdot 7\) \checkmark
\textit{Tipik hata:} annenin \(4\) yıl önceki yaşını (\(35\)) vermek.
\par\smallskip\noindent\textbf{\color{cozumyesil}Cevap: B}
\end{cozum}
\end{minipage}

\par\medskip\noindent\begin{minipage}{\linewidth}
\noindent\textbf{\color{anamavi}Soru 22.} {\small\color{gray}(Hafta 13 -- Kesir problemi; Orta)}\par
Bir öğrenci harçlığının \(\dfrac{1}{3}\)'ü ile kitap, kalan paranın \(\dfrac{1}{4}\)'ü ile defter almıştır. Geriye \(300\) TL kaldığına göre öğrencinin harçlığı kaç TL'dir?
\par\smallskip\noindent{\setlength{\tabcolsep}{1pt}\begin{tabular*}{\linewidth}{@{}l@{\extracolsep{\fill}}llll@{}}\sk{A}\(300\) & \sk{B}\(360\) & \sk{C}\(400\) & \sk{D}\(450\) & \sk{E}\(600\)\end{tabular*}}
\begin{cozum}
Kitaptan sonra kalan: \(\dfrac{2}{3}x\). Defter: \(\dfrac14\cdot\dfrac23x=\dfrac{x}{6}\). Geriye kalan:
\[\frac{2}{3}x-\frac{x}{6}=\frac{x}{2}=300\Rightarrow x=600.\]
\textit{Tipik hata:} defter harcamasını unutup \(\frac{2}{3}x=300\) yazmak \(\Rightarrow x=450\).
\par\smallskip\noindent\textbf{\color{cozumyesil}Cevap: E}
\end{cozum}
\end{minipage}

\par\medskip\noindent\begin{minipage}{\linewidth}
\noindent\textbf{\color{anamavi}Soru 23.} {\small\color{gray}(Hafta 14 -- Basit faiz; Kolay)}\par
\(40\,000\) TL, yıllık \(\%24\) basit faiz oranıyla \(9\) ay süreyle bankaya yatırılıyor. Bu süre sonunda elde edilen faiz kaç TL'dir?
\par\smallskip\noindent{\setlength{\tabcolsep}{1pt}\begin{tabular*}{\linewidth}{@{}l@{\extracolsep{\fill}}llll@{}}\sk{A}\(800\) & \sk{B}\(4800\) & \sk{C}\(7200\) & \sk{D}\(9600\) & \sk{E}\(47\,200\)\end{tabular*}}
\begin{cozum}
Süre ay cinsinden: \(F=\dfrac{A\cdot n\cdot t}{1200}=\dfrac{40\,000\cdot 24\cdot 9}{1200}=7200\) TL.
\textit{Tipik hatalar:} bir yıllık faizi (\(9600\)) vermek; faiz yerine dönem sonundaki toplam parayı (\(47\,200\)) vermek.
\par\smallskip\noindent\textbf{\color{cozumyesil}Cevap: C}
\end{cozum}
\end{minipage}

\par\medskip\noindent\begin{minipage}{\linewidth}
\noindent\textbf{\color{anamavi}Soru 24.} {\small\color{gray}(Hafta 14 -- İndirim ve kâr; Orta)}\par
Bir satıcı, maliyeti \(1250\) TL olan bir ürünün etiket fiyatını maliyetin \(\%40\) fazlası olarak belirliyor. Ürünü etiket fiyatı üzerinden \(\%20\) indirimle satan satıcının kârı kaç TL'dir?
\par\smallskip\noindent{\setlength{\tabcolsep}{1pt}\begin{tabular*}{\linewidth}{@{}l@{\extracolsep{\fill}}llll@{}}\sk{A}\(150\) & \sk{B}\(250\) & \sk{C}\(280\) & \sk{D}\(350\) & \sk{E}\(500\)\end{tabular*}}
\begin{cozum}
Etiket fiyatı: \(1250\cdot\dfrac{140}{100}=1750\) TL. İndirimli satış fiyatı: \(1750\cdot\dfrac{80}{100}=1400\) TL.
\[\text{Kâr}=1400-1250=150\text{ TL}.\]
\textit{Tipik hatalar:} yüzdeleri doğrudan çıkarıp \(\%40-\%20=\%20\) kâr sanmak \(\Rightarrow 250\); indirim tutarını (\(350\)) ya da indirimsiz kârı (\(500\)) vermek.
\par\smallskip\noindent\textbf{\color{cozumyesil}Cevap: A}
\end{cozum}
\end{minipage}

\par\medskip\noindent\begin{minipage}{\linewidth}
\noindent\textbf{\color{anamavi}Soru 25.} {\small\color{gray}(Hafta 14 -- Karışım; Zor)}\par
Tuz oranı \(\%20\) olan \(300\) gram tuzlu su ile tuz oranı \(\%35\) olan kaç gram tuzlu su karıştırılırsa tuz oranı \(\%30\) olan bir karışım elde edilir?
\par\smallskip\noindent{\setlength{\tabcolsep}{1pt}\begin{tabular*}{\linewidth}{@{}l@{\extracolsep{\fill}}llll@{}}\sk{A}\(150\) & \sk{B}\(300\) & \sk{C}\(450\) & \sk{D}\(600\) & \sk{E}\(900\)\end{tabular*}}
\begin{cozum}
Tuz miktarları korunur: \(300\cdot\dfrac{20}{100}+m\cdot\dfrac{35}{100}=(300+m)\cdot\dfrac{30}{100}\).
\(100\) ile çarpalım: \(6000+35m=9000+30m\Rightarrow 5m=3000\Rightarrow m=600\) gram.
(Kısa yol: \(\%20\), hedefin \(10\) puan altında; \(\%35\), \(5\) puan üstünde: \(300\cdot 10=m\cdot 5\Rightarrow m=600\).)
\par\smallskip\noindent\textbf{\color{cozumyesil}Cevap: D}
\end{cozum}
\end{minipage}

\end{document}
